%!TEX root = ../template.tex
% Comentarios da revisao critica de 2026-09-13, visiveis no PDF.
% Para a versao sem comentarios, trocar a linha seguinte por
% \mostrarrevisaofalse e recompilar template.tex.
\newif\ifmostrarrevisao
\mostrarrevisaotrue

\tcbuselibrary{breakable}
\newcommand{\revisaotese}[5]{%
  \ifmostrarrevisao
    \par\begingroup
    \colorlet{revisaocor}{blue!55!black}%
    \ifstrequal{#2}{P1}{\colorlet{revisaocor}{red!65!black}}{}%
    \ifstrequal{#2}{P3}{\colorlet{revisaocor}{black!60}}{}%
    \ifstrequal{#3}{FIGURA}{\colorlet{revisaocor}{teal!65!black}}{}%
    \begin{tcolorbox}[
      breakable,
      colback=revisaocor!4!white,
      colframe=revisaocor,
      colbacktitle=revisaocor,
      coltitle=white,
      boxrule=0.5pt,arc=1mm,
      left=2.5mm,right=2.5mm,top=2mm,bottom=2mm,
      before skip=8pt,after skip=8pt,
      fonttitle=\sffamily\bfseries\small,
      fontupper=\normalfont\small,
      title={REVISÃO #1 | #2 | #3}]
      \setlength{\emergencystretch}{3em}%
      \textbf{#4}\par\smallskip
      #5
    \end{tcolorbox}%
    \endgroup
  \fi
}

% As definicoes de glossarios sao lidas no preambulo. As tres notas
% sobre esses ficheiros aparecem neste apendice apenas em modo de revisao.
\ifmostrarrevisao
  % 2026-09-18: Apendice H arquivado em _revisao/ (so tinha FR-01/02/03,
  % resolvidas a 2026-09-17 com as siglas, simbolos e glossario reais).
  % \ntaddfile{appendix}{appendix-H-revisao}
\fi

% ---------------------------------------------------------------------
% 2026-09-19: marcacao temporaria das alteracoes da revisao integral.
% \alterado{...} destaca o texto novo ou reescrito. Para retirar TODAS as
% marcas de uma vez, trocar a definicao activa pela neutra, logo abaixo.
% (Nao se usa soul/\hl: rebenta com acentos, \ref e \cite em pdflatex.)
\definecolor{corAlterado}{RGB}{176,0,132}
\newcommand{\alterado}[1]{{\color{corAlterado}#1}}
%\newcommand{\alterado}[1]{#1}   %% <- versao final, sem marcas
% ---------------------------------------------------------------------
